\slajdid{10-s004}
\begin{frame}[label=10-s004]{Impulsni režim rada diode}
\fontsize{9}{10}\selectfont\setlength{\abovedisplayskip}{2pt}\setlength{\belowdisplayskip}{2pt}\setlength{\abovedisplayshortskip}{1pt}\setlength{\belowdisplayshortskip}{1pt}
\begin{columns}[c,onlytextwidth]\column{.56\textwidth}\centering\begin{tikzpicture}[x=.85cm,y=.65cm,font=\sitneoznake,>=Latex]\ctikzset{bipoles/length=.65cm}\draw(-1,1)node[ocirc]{}to[R,l=$R$](1,1)--(1,.7);\begin{scope}[shift={(1,-.05)},rotate=-90]\draw(-1,0)--(-.3,0);\draw(-.3,-.3)--(-.3,.3)--(.3,0)--cycle;\draw(.3,0)--(1,0);\draw(.3,-.3)--(.3,.3);\end{scope}\draw(1,-1.05)--(1,-1.4)node[ground]{};\draw(-1,-1.4)node[ground]{}--(-1,-.8);\draw[->](-1,-.7)--(-1,.7)node[midway,left]{$v_i(t)$}node[near end,left]{$+$};\draw(1,.65)--(2.3,.65);\draw(2.3,-1.4)node[ground]{}--(2.3,-.8);\draw[->](2.3,-.7)--(2.3,.65)node[midway,right]{$v_o(t)$}node[near end,right]{$+$};\draw[->](.35,1.4)--(.85,1.4)node[midway,above]{$i_D(t)$};\end{tikzpicture}
\hfill\input{slike/tikz/s004_ulazni_impuls.tex}\column{.42\textwidth}\centering\begin{tikzpicture}[x=.65cm,y=.30cm,font=\sitneoznake,>=Latex]
\draw[->](0,5.5)--(6.2,5.5)node[below]{t};\draw[->](.2,4.1)--(.2,7.4)node[right]{$v_i(t)$};\draw(.2,4.4)--(1,4.4)--(1,6.6)--(3.3,6.6)--(3.3,4.4)--(5.7,4.4);\node[left]at(.2,6.6){$V_H$};\node[left]at(.2,4.4){$V_L$};
\draw[->](0,2.9)--(6.2,2.9)node[below]{t};\draw[->](.2,1.5)--(.2,4.1)node[right]{$v_o(t)$};\draw[dashed](.2,3.4)--(3.3,3.4);\draw[dashed](.2,1.8)--(5.7,1.8);\node[left]at(.2,3.4){$V_D$};\node[left]at(.2,1.8){$V_L$};
\draw(1,1.8)..controls(1.2,3.25)and(1.45,3.4)..(2,3.4)--(3.3,3.4)--(4.8,3.2)--(4.8,2.95)..controls(5,2.15)and(5.6,1.85)..(6,1.8);
\draw[->](0,-.1)--(6.2,-.1)node[below]{t};\draw[->](.2,-1.7)--(.2,1.4)node[right]{$i_D(t)$};\draw(1,.95)--(1.6,.95)..controls(1.8,.95)and(1.8,.8)..(2.2,.8)--(3.3,.8)--(3.3,-.95)--(4.8,-.75)..controls(5,-.3)and(5.6,-.15)..(6,-.1);
\foreach\x/\n in{1/0,1.6/1,3.3/2,4.8/3,6/4}{\draw[dashed](\x,-1.5)--(\x,4.3);\node[below]at(\x,-1.5){$t_\n$};}
\end{tikzpicture}
\end{columns}\smallskip
Period $t_0$ do $t_1$ – vreme kašnjenja; puni se oblast prostornog tovara, „kapacitivnost“, koja je bila inverzno polarizovana i „skuplja“ se, još uvek je inverzno polarizovan spoj.
\par
Period $t_1$ do $t_2$ – pri kraju prethodnog perioda pn spoj je direktno polarizovan. Dioda vodi. Struja je određena veličinom ulaznog napona, otpornošću i naponom direktne polarizacije diode
 $I_D=\frac{V_H-V_D}{R}$\par
Period $t_2$ do $t_3$ – vreme nagomilavanja. Postoji velika koncentracija nagomilanih nosilaca uz oblast prostornog tovara. Spoj je i dalje direktno polarizovan. Pojavljuje se negativna struja koja odvodi nagomilane nosioce.
 $I_D=\frac{V_L-V_D}{R}\quad i\quad V_L<0\ \to\ I_D<0$\par
Period $t_3$ do $t_4$ – vreme prelaza. Oblast pn spoja postaje inverzno polarizovano i širi se. Prazne se oblasti prostornog tovara „kapacitivnosti“. Zato i jeste eksponencijalan oblik.
\end{frame}
